% 原PDF181—183页全部14道习题；184页为出版社标识页。
\clearpage
\section*{习题}
\addcontentsline{toc}{section}{习题}
\markright{习题}

\begin{enumerate}[label=\textbf{习题\arabic*.},leftmargin=5em]
\item 考虑例4.1的原始计数数据。假设 $P=0.75$，使用上分位数标准化计算标准化因子和标准化后的计数。

\item 证明
\[
\operatorname{var}(M_g(j,r))=\frac{D_j-n_{gj}}{D_jn_{gj}}+
\frac{D_r-n_{gr}}{D_rn_{gr}}.
\]
\item 考虑例4.1的原始计数数据。使用M值截尾均值计算标准化因子和标准化后的计数。

\item 证明优化问题（4.1A49）
\[
\begin{aligned}
\max_{q_\theta(\theta)}\quad&F(q_z^{(u+1)}(Z),q_\theta(\theta))\\
\text{约束为}\quad&\int q_\theta(\theta)\,d\theta=1
\end{aligned}
\]
的解为
\[
q_\theta^{(u+1)}(\theta)=\frac1{C_\theta}P(\theta^{(u)})
\exp\left[\int q_z(Z^{(u+1)})\log P(Z,R,Q\mid\theta^{(u)})\,dZ\right],
\]
其中
\[
C_\theta=\int P(\theta^{(u)})\exp\left[
\int q_z(Z^{(u+1)})\log P(Z,R,Q\mid\theta^{(u)})\,dZ\right].
\]
\item 证明
\[
\E_\theta[\log\theta_i]=\psi(\alpha_i)-\psi\left(\sum_{j=1}^{M}\alpha_j\right),
\qquad\psi(\alpha)=\frac{d\Gamma(\alpha)/d\alpha}{\Gamma(\alpha)},
\]
其中 $\theta$ 服从Dirichlet分布。

\item 证明Dirichlet分布（式（4.A64））的均值为
\[
\E_\theta[\theta_i^{(u+1)}]=\frac{\alpha_i^{(u+1)}}{\sum_{i=0}^{M}\alpha_i^{(u+1)}}.
\]
\item 证明结论4.4，即有缺口比对RNA测序数据的变分贝叶斯算法。

\item 证明Beta分布的均值为
\[
\E[\phi]=\frac{\beta_{i1}}{\beta_{i1}+\beta_{i2}}.
\]
\item 证明具有泊松分布的对数线性模型，其对数似然取得极大值的必要条件为
\[
\frac{\partial l(N_{ig}\mid\alpha_g)}{\partial\alpha_g}
=\sum_i[N_{ig}y_i-N_{ig}^{(0)}y_i\exp(y_i\alpha_g)]=0,
\]
\[
\begin{aligned}
\frac{\partial l(N_{ig}\mid\alpha_{gm})}{\partial\alpha_{gm}}
&=\sum_i\left[N_{ig}I_{(i\in C_m)}-N_{ig}^{(0)}I_{(i\in C_m)}
\exp\left(\sum_m\alpha_{gm}I_{(i\in C_m)}\right)\right]\\
&=0.
\end{aligned}
\]
\item 证明
\[
\E[U(\theta)]=0,
\]
以及
\[
\operatorname{var}(U(\theta))=-\E\left[
\frac{\partial^2U}{\partial\theta\,\partial\theta^T}\right]=I(\theta).
\]
\item 证明检验基因 $g$ 与二分类或多分类结局关联的得分统计量为
\[
S_{cg}=\sum_{m=1}^{M}\frac{
\left[\sum_{i\in C_m}(N_{ig}-N_{ig}^{(0)})\right]^2}{
\sum_{i\in C_m}N_{ig}^{(0)}}.
\]
\item 证明：若 $y_{gi}\sim\mathrm{NB}(r_i,p)$，则
$n_{gA}\sim\mathrm{NB}(\sum_{i=1}^{m_A}r_i,p)$。

\item 证明
\[
\|U\Lambda V^T-U\Sigma V^T\|_F^2=\|\Lambda-\Sigma\|_F^2.
\]
\item 证明以下两个优化问题等价：
\[
\begin{aligned}
L(\delta_i,b^i,Z)={}&\frac12f(\delta_i,b^i)\\
&{}+\lambda_1\left(\sum_{j=1}^{i-1}\|Z_j^1\|_2+\sum_{j=i+1}^{K}\|Z_j^1\|_2\right)\\
&{}+\lambda_2\sum_{l=1}^{r_i}\|Z_l^2\|_2
+\sum_{j=1}^{i-1}\mu_j^T(\delta_{ji}-Z_j^1)\\
&{}+\sum_{j=i+1}^{K}\mu_j^T(\delta_{ji}-Z_j^1)\\
&{}+\frac{\rho_1}{2}\left(\sum_{j=1}^{i-1}\|\delta_{ji}-Z_j^1\|_2^2
+\sum_{j=i+1}^{K}\|\delta_{ji}-Z_j^1\|_2^2\right)\\
&{}+\sum_{l=1}^{r_i}\pi_l^T(b_l^i-Z_l^2)
+\frac{\rho_2}{2}\sum_{l=1}^{r_i}\|b_l^i-Z_l^2\|_2^2,
\end{aligned}
\]
以及
\[
\begin{aligned}
L(\delta_i,b^i,Z)={}&\frac12f(\delta_i,b^i)\\
&{}+\lambda_1\left(\sum_{j=1}^{i-1}\|Z_j^1\|_2+\sum_{j=i+1}^{K}\|Z_j^1\|_2\right)
+\lambda_2\sum_{l=1}^{r_i}\|Z_l^2\|_2\\
&{}+\frac{\rho_1}{2}\left(\sum_{j=1}^{i-1}\|\delta_{ji}-Z_j^1+u_j\|_2^2
+\sum_{j=i+1}^{K}\|\delta_{ji}-Z_j^1+u_j\|_2^2\right)\\
&{}+\frac{\rho_2}{2}\sum_{l=1}^{r_i}\|b_l^i-Z_l^2+v_l\|_2^2.
\end{aligned}
\]
\end{enumerate}
